\section{Appendix: Explainability diagnostics}
\label{app:explainability-diagnostics}

This appendix summarizes the calculations used to construct the Integrated Gradients (IG) and decoder-layer PCA diagnostics. Let \(X\) denote one normalized model input sequence and \(X^{0}\) its IG baseline. In this analysis the default baseline is \(X^{0}=0\), which corresponds to the training-mean input in normalized space. The model prediction is denoted by \(\hat{q}=f(X)\), where \(\hat{q}_{y,x}\) is the predicted submesoscale vertical heat flux at output grid cell \((y,x)\). For each selected snapshot, three output cells are explained: the strongest positive flux, the strongest negative flux, and the weakest absolute flux. Let \(m=1,\ldots,M\) index all explained target cells, where \(M=N_{\mathrm{snap}}\times 3\). In the present calculation, \(N_{\mathrm{snap}}=90\), so \(M=270\).

For target cell \(m\), the scalar objective is
\[
J_m(X)=\hat{q}_{y_m,x_m},
\]
where \((y_m,x_m)\) is the explained target location. Let \(\ell=1,\ldots,L\) index temporal offsets, \(c=1,\ldots,C\) index input features, and \((j,i)\) index input grid points in the meridional and zonal directions. The eight features are \(U\), \(V\), \(W\), vorticity, strain, potential temperature, horizontal speed, and KPP boundary-layer depth. With \(K=64\) integration steps and \(\alpha_k=k/K\), the IG attribution is computed as
\[
A^{(m)}_{\ell,c,j,i}
=
\left(X_{\ell,c,j,i}-X^{0}_{\ell,c,j,i}\right)
\frac{1}{K}\sum_{k=1}^{K}
\frac{\partial J_m\!\left(X^{0}+\alpha_k(X-X^{0})\right)}
{\partial X_{\ell,c,j,i}} .
\]
Positive and negative values of \(A^{(m)}_{\ell,c,j,i}\) indicate input components that increase or decrease the selected scalar prediction, respectively.

\subsection{Absolute IGs for geographical feature attribution}
\label{app:absolute-igs-geographical}

The geographical attribution map measures where each input feature contributes strongly in the fixed model domain, independent of the sign of the contribution. For feature \(c\), the plotted geographical field is
\[
G_{c,j,i}
=
\frac{1}{M}\sum_{m=1}^{M}\sum_{\ell=1}^{L}
\left|A^{(m)}_{\ell,c,j,i}\right| .
\]
Here \(G_{c,j,i}\) is the mean absolute IG contribution of feature \(c\) at physical grid point \((j,i)\), averaged over all explained target cells and all temporal offsets. The absolute value is used because this panel is intended to quantify attribution magnitude rather than whether the feature increases or decreases the predicted heat flux.

\subsection{Signed IGs at explained target cells}
\label{app:signed-igs-targets}

The signed IG map keeps the sign of the contribution and assigns the total signed feature contribution back to the explained output cell. For target \(m\), the signed contribution of feature \(c\) is
\[
s^{(m)}_{c}
=
\sum_{\ell=1}^{L}\sum_{j,i} A^{(m)}_{\ell,c,j,i}.
\]
Let \(n_{j,i}\) be the number of explained targets located at output cell \((j,i)\). The signed map plotted for feature \(c\) is
\[
S_{c,j,i}
=
\frac{1}{n_{j,i}}
\sum_{m:\,(y_m,x_m)=(j,i)} s^{(m)}_{c},
\]
for grid cells with \(n_{j,i}>0\). Thus \(S_{c,j,i}>0\) means that feature \(c\), summed over all input lags and input grid points, tends to increase the predicted flux at explained targets located at \((j,i)\), while \(S_{c,j,i}<0\) indicates a decreasing contribution.

\subsection{Spatial locality from target-centered IGs}
\label{app:spatial-locality-target-centered}

To assess spatial locality, each absolute attribution map is recentered around the explained prediction cell before averaging. First, for target \(m\) and feature \(c\), define the lag-summed absolute attribution
\[
F^{(m)}_{c,j,i}
=
\sum_{\ell=1}^{L}\left|A^{(m)}_{\ell,c,j,i}\right| .
\]
The map \(F^{(m)}_{c,j,i}\) is shifted so that the target cell \((y_m,x_m)\) lies at the origin \((\Delta y,\Delta x)=(0,0)\). The zonal direction is treated as periodic when shifting, while the meridional direction is not; meridional points shifted outside the domain are treated as missing. The target-centered map is then
\[
C_{c,\Delta y,\Delta x}
=
\frac{1}{N_{\Delta y,\Delta x}}
\sum_{m=1}^{M}
\widetilde{F}^{(m)}_{c,\Delta y,\Delta x},
\]
where \(\widetilde{F}^{(m)}\) is the recentered field and \(N_{\Delta y,\Delta x}\) is the number of valid targets contributing to offset \((\Delta y,\Delta x)\).

For the spatial-locality radius, let
\[
r(\Delta y,\Delta x)
=
\sqrt{(\Delta y\,\Delta_{\mathrm{km}}^{y})^2
+(\Delta x\,\Delta_{\mathrm{km}}^{x})^2},
\]
where \(\Delta_{\mathrm{km}}^{x}=\Delta_{\mathrm{km}}^{y}=50\) km in the plotted diagnostics. For feature \(c\), the cumulative attribution within radius \(r_0\) is
\[
P_c(r_0)
=
\frac{
\sum_{\Delta y,\Delta x:\,r(\Delta y,\Delta x)\le r_0}
C_{c,\Delta y,\Delta x}
}{
\sum_{\Delta y,\Delta x} C_{c,\Delta y,\Delta x}
}.
\]
The reported \(r_{50}\) and \(r_{90}\) are the smallest radii satisfying \(P_c(r_{50})\ge 0.5\) and \(P_c(r_{90})\ge 0.9\), respectively.

\subsection{Time-lag IGs}
\label{app:time-lag-igs}

The time-lag diagnostic measures which input times carry the largest absolute attribution. The feature-resolved lag attribution is
\[
L_{\ell,c}
=
\frac{1}{M}\sum_{m=1}^{M}\sum_{j,i}
\left|A^{(m)}_{\ell,c,j,i}\right| .
\]
The left panel of the time-lag figure plots the fraction of total attribution assigned to lag \(\ell\):
\[
p_{\ell}
=
\frac{\sum_{c=1}^{C}L_{\ell,c}}
{\sum_{\ell'=1}^{L}\sum_{c=1}^{C}L_{\ell',c}} .
\]
The heatmap uses a within-feature normalization,
\[
Q_{\ell,c}
=
\frac{L_{\ell,c}}
{\sum_{\ell'=1}^{L}L_{\ell',c}},
\]
so that each feature column shows how that feature's attribution is distributed across temporal offsets.

\subsection{PCA of the final decoder layer}
\label{app:decoder-pca}

Decoder-layer PCA was applied to the final decoder feature output from the trained encoder--ConvLSTM model. Let \(Z\in\mathbb{R}^{N\times C_d\times H_d\times W_d}\) denote the captured decoder tensor, where \(N\) is the number of samples, \(C_d\) is the number of decoder channels, and \((H_d,W_d)\) are the decoder-map dimensions. In the plotted example, the tensor is captured from the final decoder block, denoted \texttt{up1}, for test sample 9523.

The decoder tensor is reshaped into a matrix
\[
Z^{\ast}\in\mathbb{R}^{(N H_d W_d)\times C_d},
\]
so that each spatial location is one observation and each decoder channel is one variable. The channel mean
\[
\mu_c=\frac{1}{N H_d W_d}\sum_{n,j,i} Z_{n,c,j,i}
\]
is subtracted to form the centered matrix \(\bar{Z}^{\ast}\). Singular value decomposition is then applied:
\[
\bar{Z}^{\ast}=U\Sigma V^{\mathsf{T}}.
\]
The principal-component score for mode \(k\) at each spatial location is given by the \(k\)-th column of \(U\Sigma\), reshaped back to \((N,H_d,W_d)\). These reshaped score maps are the plotted decoder PCA modes. The explained variance ratio of mode \(k\) is
\[
\rho_k
=
\frac{\sigma_k^2/(N H_d W_d-1)}
{\sum_{r}\sigma_r^2/(N H_d W_d-1)}
=
\frac{\sigma_k^2}{\sum_{r}\sigma_r^2},
\]
where \(\sigma_k\) is the \(k\)-th singular value. Correlations with physical diagnostics are computed as spatial Pearson correlations between each PC score map and each diagnostic field after flattening valid grid cells.
